\section{Memory with Privacy}
\sectionkicker{PRIVATE COMPUTE}
\begin{wideflow}
{\Large\bfseries 覚える雲、鍵は手元に}\par
\smallskip
{\color{SurveyMuted}9月23日の私的計算の構想を、提供の有無と切り分けて示す。}\par
\end{wideflow}

9月23日、Googleは私的計算に越し方の記憶を持たせる構想を示した\cite{googlemem}。鍵は利用者の手元に置き、雲の隔離の座（enclave）でだけ解いて用を足し、すぐ包み直す。使った道具の正しさを手元で確かめられる公開の記録も添えられ、外の監査の結果にも触れられる。DeepMindと端末・中核・雲の四つの組の作とされる。これまでの私的計算が用を足すたびに忘れる仕立てだったのに対し、跨いで覚える道を開く構想である。

ただし書きぶりは「これからできる」の構想であり、今日使える品の知らせではない。窓を過ぎた場の声にも「まだ構想」と読む向きがある。監査の担い手や中身はこの紙からは確かめられず、技の書（technical brief）も表紙の域に留める。できることの広がりも、使える日の目処も、本号では言わない。

\begin{claimboundary}[資料と評価の境界]
構想の表明であり、提供の事実ではない。守りの性質は作り手の述べるところであり、他所の確かめを待つ。使える日の目処やできることの広がりは断定しない。
\end{claimboundary}
